% Auto-generated by src/make_tables.py -- do not edit by hand.
\begin{table}[t]
\centering
\caption{Order against monotonicity, on functions whose gradients are known. \emph{cubic} is the unlimited reconstruction of Proposition~1 with a probe far below the mesh; \emph{limited} is the same gradient passed through a monotonised-central limiter built from the nodal one-sided slopes. \emph{Wrong way} is the fraction of samples where the reconstruction has the OPPOSITE sign to a significant true gradient --- for a control, which is a derivative, that is a trade in the wrong direction rather than a small error, and returning zero at a turning point is the safe answer and is not counted. \emph{Overshoot} is the largest recovered gradient relative to the largest true one. The order column is given only where an order exists: a discontinuous gradient and noisy data have none. On smooth data the limiter is inert (no wrong-way samples either way) and costs nothing except at interior turning points, where clipping drops it from $3.7$ to $1.5$. On data that are not smooth it is what stops the control pointing the wrong way.}
\label{tab:limit}
\scriptsize
\setlength{\tabcolsep}{3pt}
\begin{tabular}{lcccc}
\toprule
 & \multicolumn{2}{c}{wrong way (\%)} & \multicolumn{2}{c}{overshoot ($\times$)} \\
\cmidrule(lr){2-3} \cmidrule(lr){4-5}
data condition & cubic & limited & cubic & limited \\
\midrule
smooth, monotone & 0.0 & 0.0 & 1.00 & 1.00 \\
smooth, interior turning point & 0.0 & 0.0 & 1.01 & 1.00 \\
gradient kink & 0.2 & 0.0 & 1.33 & 1.24 \\
2\% nodal noise (solved grid) & 8.6 & 0.0 & 2.26 & 1.64 \\
\midrule
\multicolumn{5}{l}{\emph{convergence order, where one exists}} \\
\multicolumn{5}{l}{smooth, monotone: cubic $3.41$, limited $3.41$} \\
\multicolumn{5}{l}{smooth, interior turning point: cubic $3.72$, limited $1.51$} \\
\bottomrule
\end{tabular}
\end{table}
